\documentclass[11pt,letterpaper]{article}
\usepackage[margin=1in,headheight=14pt,headsep=18pt]{geometry}
\usepackage{fontspec}
\usepackage[UTF8,scheme=plain,fontset=fandol]{ctex}
\setmainfont{texgyretermes}[Extension=.otf,UprightFont=*-regular,BoldFont=*-bold,ItalicFont=*-italic,BoldItalicFont=*-bolditalic]
\usepackage{setspace}
\usepackage{graphicx}
\usepackage{longtable,booktabs,array}
\usepackage{calc}
\usepackage{amsmath,amssymb}
\usepackage{fancyhdr}
\usepackage{titlesec}
\usepackage{needspace}
\usepackage{flafter}
\usepackage{placeins}
\usepackage{xurl}
\usepackage[unicode,hidelinks,breaklinks]{hyperref}
\urlstyle{same}
\hypersetup{pdfauthor={},pdftitle={基于机器学习的微博用户主观幸福感评估工具开发及其心理学意义：一项评价性与解释性研究},pdfsubject={Journal-style computational reanalysis draft}}
\pagestyle{fancy}
\fancyhf{}
\fancyfoot[R]{\thepage}
\renewcommand{\headrulewidth}{0pt}
\setlength{\parindent}{2em}
\setlength{\parskip}{3pt plus 1pt}
\setlength{\emergencystretch}{3em}
\setcounter{secnumdepth}{0}
\titleformat{\section}{\normalfont\large\bfseries}{}{0pt}{}
\titleformat{\subsection}{\normalfont\normalsize\bfseries}{}{0pt}{}
\titleformat{\subsubsection}{\normalfont\normalsize\bfseries\itshape}{}{0pt}{}
\titlespacing*{\section}{0pt}{13pt}{5pt}
\titlespacing*{\subsection}{0pt}{10pt}{3pt}
\titlespacing*{\subsubsection}{0pt}{8pt}{3pt}
\providecommand{\tightlist}{\setlength{\itemsep}{0pt}\setlength{\parskip}{0pt}}
\providecommand{\pandocbounded}[1]{#1}
\newcommand{\APAref}{\par\noindent\hangindent=0.5in\hangafter=1\relax}
\setlength{\LTpre}{4pt}
\setlength{\LTpost}{4pt}
\setlength{\tabcolsep}{4pt}
\setlength{\textfloatsep}{12pt plus 2pt minus 2pt}
\setlength{\intextsep}{10pt plus 2pt minus 2pt}
\renewcommand{\topfraction}{0.85}
\renewcommand{\bottomfraction}{0.75}
\renewcommand{\textfraction}{0.12}
\renewcommand{\floatpagefraction}{0.72}
\setstretch{1.30}
\widowpenalty=10000
\clubpenalty=10000
\begin{document}
\begin{center}\fontsize{16}{20}\selectfont\bfseries 基于机器学习的微博用户主观幸福感评估工具开发及其心理学意义：一项评价性与解释性研究\end{center}
\section{摘要}\label{ux6458ux8981}

社交媒体语言可能包含与主观幸福感有关的信息，但预测关联、数值准确性与测量效度需要不同证据。我们对 1,037 名微博用户开展二次分析，使用 120 个心理语言特征估计消极情绪、积极情绪和生活满意度。在固定搜索预算内，采用五折交叉验证比较八类模型，并保留 104 名参与者用于评估。选定的共享随机森林在留出集上的相关分别为 0.247、0.284 和 0.152，均方根误差分别为 0.166、0.165 和 0.202，单位为数据保存的分数单位。两种情绪的相关在校正后仍显著，而生活满意度的区间包含零。相对训练均值基线的误差降幅较小，未达到多重比较校正标准；共享学习相对独立森林也未显示明确优势。随后开展的探索性搜索未改善嵌套开发集表现。归因排序识别出模型使用的类别，但相应的留出集置换效应较小且方向混合。预定敏感性分析仍呈现有限的预测表现。这些结果支持适度的回顾性语言---情绪关联，但尚未确立准确的个体评估、时间信度或可迁移性。

关键词：主观幸福感；心理语言学；社交媒体；机器学习；多输出回归；可复现性

利益冲突：待作者提供声明。

\section{引言}\label{ux5f15ux8a00}

主观幸福感关注人们如何体验和评价自己的生活。其情感成分与认知成分相互关联，但又有所区别：积极情绪、消极情绪和对生活满意度的总体判断分别反映主观体验的不同方面 (Diener, 1984)。积极与消极情绪量表（PANAS）分别测量积极情绪和消极情绪，生活满意度量表（SWLS）则评估对生活满意程度的总体认知判断 (Diener et al., 1985; Watson et al., 1988)。因此，基于语言的幸福感评估应当分别检验各个成分，而不能假定对某一成分的成功预测即可确立对全部三个成分的效度。

这种分维度的思路也界定了拟议工具的测量范围。Diener 等（1999）强调，幸福感各成分可能具有不同的相关因素；Ryan 和 Deci（2001）则区分享乐取向的评价与意义、功能等实现取向的关切。因此，即使能够预测给定的三个分数，也不能据此认定已经全面测量了蓬勃发展、德性或心理功能。相关效标仍是参与者如何报告自己的情绪、评价自己的生活，而非根据其谈论的话题推断一种规范性判断（Diener et al., 2018）。

社交媒体语言提供了可能包含上述成分相关信息的行为记录。Schwartz et al. (2013) 发现 Facebook 语言与人格及人口学特征之间存在关联；Park et al. (2015) 则从聚合效度、区分效度、知情者报告、外部效标和时间证据等方面，评估了基于语言的人格分数。Eichstaedt et al. (2018) 考察了另一种应用，即通过 Facebook 语言预测后续医疗记录中的抑郁记录。这些研究共同为考察语言中的心理信息提供了依据。同时，它们也表明，预测证据依赖具体的构念、效标、人群与时间设计；在这些情境中的表现，并不能直接确立从微博语言推导出的主观幸福感分数的效度。

心理语言词典提供了一种可解释的语言表征。预定义类别把词语使用与可识别的心理领域相联系，同时也会丢失部分语境信息。Tausczik 和 Pennebaker（2010）说明，语境、反讽和习语可能使简单的词语类别归属产生偏差。封闭词汇与开放词汇方法的比较进一步表明，宽泛类别概括与更具体的内容模式可以承担互补作用（Eichstaedt et al., 2021）。在中文社交媒体文本中，Zhao 等（2016）发现，简体中文 LIWC 类别分数与人工判断的一致程度因类别而异。因此，检出与某种情绪相关的词语，并不等于测量了作者的情绪体验。基于词典的预测既需要对照参与者层面的效标进行评估，也需要检查被概括的语言内容。

主观幸福感的多维结构还引出如何建模这些特征的问题。通过共享表征学习相关结局能够提供归纳偏置，这一思路在多任务学习中得到发展（Caruana, 1997）。在本研究中，独立回归器允许每个成分具有不同的语言关系，而共享输出树使用共同的特征空间划分并返回结局向量（Kocev et al., 2007）。我们将其具体称为多输出回归：更广泛的多任务学习文献还涵盖超出这一情境的共享结构和任务定义（Zhang \& Yang, 2022）。结局之间的相关性构成比较的动机，却不能证明共同划分能够获得更好的泛化表现。因此，模型比较必须使用未参与拟合和参数选择的参与者（Yarkoni \& Westfall, 2017）。

本研究通过对既有微博数据集的二次分析考察这些问题。我们关注：聚合语言特征能否预测留出参与者的各个幸福感成分；在固定搜索预算内，选定的共享输出程序能否优于独立程序；以及哪些类别对选定树模型的预测作出了贡献。这些问题区分了预测关联、数值误差改善和对拟合模型的解释。按参与者划分数据并在训练折内拟合预处理，可以应对重复观察和信息泄漏带来的评估风险 (Kapoor \& Narayanan, 2023)；与均值基线比较并报告不确定性，则有助于判断所得证据的强度。

\section{方法}\label{ux65b9ux6cd5}

\subsection{参与者与数据来源}\label{ux53c2ux4e0eux8005ux4e0eux6570ux636eux6765ux6e90}

随附研究说明记载，研究通过向约 20,000 名微博用户发送私信邀请进行便利招募。知情同意的参与者完成一次在线问卷，随后研究者通过新浪微博 API 获取其公开帖子和个人资料。说明中的纳入条件排除了未满 18 岁、没有公开账号或未提供同意，以及问卷作答中断的用户；完成者获得了金钱奖励。数据采集与一个于 2012 年九月至 2013 年二月开展的平行项目有关，但无法独立确认本文件对应的确切招募时段。上述程序据数据采集说明报告；本研究使用既有数据，未招募新的参与者。

原文件包含 1,427 行数据，对应 1,037 个不同的 ID。其中 195 个 ID 各有两条不同的观察记录，并有其中一条记录的完全相同副本，由此形成 195 行重复记录。在模型比较前，我们删除精确重复记录，并为每名参与者保留数值型问卷时间戳最早的观察。不同记录在该最早时间戳上不存在并列。处理后得到 1,037 名参与者，避免了因重复观察而赋予某些个体更高权重。较晚记录保留用于敏感性分析。人口学变量和问卷题目不可用，也无法重建招募各阶段的人数。

原始采集说明记载了知情同意、自愿参与、退出权利及机构伦理审查批准。可识别机构的批准信息保留在独立标题页中。相关批准和同意文件、此次二次使用的授权，以及数据发布权限，仍需由负责的作者确认。

\subsection{测量与语言特征}\label{ux6d4bux91cfux4e0eux8bedux8a00ux7279ux5f81}

数据字典将 SWB\_\allowbreak{}nn 定义为消极情绪，SWB\_\allowbreak{}pp 定义为积极情绪，SWB\_\allowbreak{}avg 定义为生活满意度，来源于 PANAS 和 SWLS 的中文改编版本。原量表包括两个各含十个条目的情绪分量表，以及一个五条目的生活满意度量表 (Diener et al., 1985; Watson et al., 1988)。数据字典注明了一项 SWLS 中文翻译来源，但实际施测版本、情绪报告时间窗和计分变换均未得到独立核实。三个已提供分数的范围均为 0--1。我们保留这些存储单位，而未推测其与问卷总分之间的换算关系；因此，RMSE 尚不能以原始 PANAS 或 SWLS 单位作出经核实的解释。

预测变量由 AGENCY 起至最后一列的 120 个字段构成。根据采集说明，这些特征使用简体中文 LIWC、道德基础、道德动机、文化价值和基本情绪词典，汇总了各参与者聚合帖子中的类别频率。ID、时间戳、发帖量、问卷作答时长及其他问卷分数均未作为模型输入。分析从聚合特征开始，未重新进行文本提取。词典的确切版本及 13 个缩写代码的映射关系不可用。34 名参与者的 ind 特征值超过一，与简单比例的解释不符；主分析保留该特征，并在预定敏感性分析中将其与未映射代码一并排除。预测变量和结局均无缺失值或非有限值，也没有常量预测变量。未删除任何统计离群值。

在保留的 1,037 条记录中，有 1,027 条的最新帖子时间戳晚于问卷时间戳。由于原始帖子不可用，无法重建截止于问卷之前的文本时间窗。因此，本分析估计的是回顾性的参与者间关联。词级频率、奇数与偶数帖子特征以及题目作答也不可用，因而无法生成词云、计算文本分半信度或问卷内部一致性。

\subsection{模型开发}\label{ux6a21ux578bux5f00ux53d1}

分析计划在完成数据结构审计后、模型比较前固定，但未进行公开预注册。我们使用随机种子 20260923，按参与者随机分配 90\% 的数据至开发集（n = 933），其余 10\% 留作最终测试集（n = 104）。开发集内的五折交叉验证用于参数选择。所有模型和敏感性分析均沿用同一参与者划分（图 1）。

\begin{figure}[!htbp]\centering
\includegraphics[width=\linewidth,height=0.54\textheight,keepaspectratio]{figures/figure1_analysis_flow.pdf}
\par\vspace{5pt}{\raggedright\singlespacing\fontsize{10}{12}\selectfont\noindent 图 1. 参与者层面的分析流程。记录核对处理后，1,037 名参与者各贡献一条最早观察。开发集（n = 933）与留出测试集（n = 104）的 ID 不重叠。需要拟合的预处理仅在训练折中进行，模型选择先于最终评估。\par}
\end{figure}

\FloatBarrier

中位数插补器和最小---最大输入缩放器均在各训练折内拟合，再应用于相应验证折。由于数据完整，插补器未改变本数据集。结局使用训练折的均值和标准差进行标准化，使三个目标在共享输出损失中具有可比权重；预测值随后转换回存储分数单位。该折内处理避免在验证参与者上拟合预处理，也不同于在交叉验证前仅进行一次归一化。

我们比较了八类模型：训练集均值基线、独立岭回归、独立随机森林、独立极端随机树、独立直方图梯度提升、共享输出随机森林、共享输出极端随机树，以及多任务弹性网。随机森林对随机化树的预测进行平均，而极端随机树对划分阈值采用更强的随机化 (Breiman, 2001; Geurts et al., 2006)。共享树模型对所有结局使用共同的划分；独立模型则分别拟合回归器。多任务弹性网提供了另一种共享稀疏结构的线性模型。

固定参数网格包含 42 种配置和 210 次复合折内评估，每次独立目标拟合包含三个回归器。选择标准是使交叉验证归一化 RMSE 的平均值最小：先将各目标的 RMSE 除以其训练折标准差，再对结局和折求平均。独立模型为各结局分别选择参数，共享模型则依据宏平均选择一组设置。我们在不查看留出测试集表现的情况下，保留了整体选定模型、最佳独立集成模型、最佳共享模型及最佳树模型。完整搜索网格、实现检查和选定设置见补充材料。优化在固定预算耗尽后结束，未根据测试结果重新调参。

\subsection{原始评估后的探索性优化}\label{ux539fux59cbux8bc4ux4f30ux540eux7684ux63a2ux7d22ux6027ux4f18ux5316}

查看原始留出结果后，我们另行固定了包含 22 个配置的搜索，用以考察更强正则化、平滑非线性回归及低维替代方案。候选包括均值与原共享森林参照、使用原始或对数变换特征的岭回归、支持向量回归、浅层提升、弹性网、偏最小二乘，以及叶节点更大的独立森林。所有需要学习的预处理仍限于各训练折。原开发集内的五个外层折各自重复内层五折选择；随后在全部开发参与者上进行最终五折搜索，冻结一个候选后，才在此前已查看的留出集上评估。在外层折内重新执行模型选择，可以处理一类特定的选择偏差（Cawley \& Talbot, 2010）。但由于此前的开发集分析和留出结果已经知晓，整个扩展仍属于探索性分析。我们保留原主分析结果，并在固定预算用尽后停止。补充材料报告完整配置网格、外层折结果和探索性比较。

\subsection{评估与不确定性}\label{ux8bc4ux4f30ux4e0eux4e0dux786eux5b9aux6027}

模型选择完成后，所有模型类别均在同一留出测试集上评估。Pearson r 用于量化线性关联，并报告基于 Fisher-z 的 95\% 置信区间和双侧检验。RMSE 用于量化存储分数单位下的数值预测误差，并通过 5,000 次参与者 bootstrap 重抽样获得百分位数法 95\% 区间。选定模型的三个相关检验采用 Holm 校正。此外还保留了决定系数和 Spearman 相关。常量均值预测的相关系数无定义；相关本身也不能证明数值一致性或对个体内变化的敏感性。

两项预定比较分别检验选定模型与训练均值，以及最佳共享模型与最佳独立集成模型的差异。我们使用 5,000 次参与者 bootstrap 重抽样估计配对 RMSE 差值，并进行 10,000 次参与者内预测值交换，以检验均方误差差异，同时采用加一校正。Holm 校正覆盖两类比较与三个结局形成的六项检验。Bootstrap 区间是逐项区间，而非覆盖这些检验的同时区间。推断以已拟合模型和本次参与者划分为条件；交换检验还依赖相应的预测可交换性假设。这两类程序均未纳入在其他样本上重新训练所产生的不确定性，也未涵盖便利招募带来的不确定性。

\subsection{敏感性分析与模型解释}\label{ux654fux611fux6027ux5206ux6790ux4e0eux6a21ux578bux89e3ux91ca}

三项预定敏感性分析保留了原参与者分配和选定超参数，分别使用每名参与者最新的不同记录、仅保留出现过一次的 ID，以及排除 ind 和 13 个未映射代码。各模型重新拟合，但不重新调参。另进行了三次重复五折开发集交叉验证，以描述固定设置下的变异。由于开发集已用于参数选择，这些重复分析属于诊断，而不是独立的嵌套验证。

我们在开发集参与者上，使用树路径依赖型 TreeSHAP 解释开发阶段选定的最佳树模型 (Lundberg et al., 2020)。贡献值被转换回存储结局单位，核查其能否加和重建预测，并按平均绝对值排序。随后，将每个结局排名前十的特征取并集，在留出测试集中逐个置换，每个特征洗牌 20 次。洗牌变异描述的是模型敏感性，而非总体置信区间。归因和置换结果针对已拟合的预测函数，不能识别语言对幸福感的因果效应 (Janzing et al., 2020)。

分析使用 Python 3.12.7、scikit-learn 1.9.1 和 SHAP 0.52.0。脚本、参与者划分、预测值及计算复现结果均保存在配套分析包和补充材料中。

\FloatBarrier

\section{结果}\label{ux7ed3ux679c}

\subsection{样本特征与模型选择}\label{ux6837ux672cux7279ux5f81ux4e0eux6a21ux578bux9009ux62e9}

完整样本、开发集和留出测试集的结局分布见表 1。消极情绪与积极情绪的相关为 \mbox{-0.307}，与生活满意度的相关为 \mbox{-0.408}；积极情绪与生活满意度的相关为 0.458。这些关联为联合评估多个成分提供了动机，同时仍需保留各结局的独立表现估计。

\par\medskip\noindent\begin{minipage}{\linewidth}\singlespacing\fontsize{10}{12}\selectfont\setlength{\parindent}{0pt}
{\bfseries 表 1 参与者数据集中的结局分布}\par\vspace{5pt}
\begin{tabular}{@{}
  >{\raggedright\arraybackslash}p{(\linewidth - 6\tabcolsep) * \real{0.2500}}
  >{\raggedright\arraybackslash}p{(\linewidth - 6\tabcolsep) * \real{0.2500}}
  >{\raggedright\arraybackslash}p{(\linewidth - 6\tabcolsep) * \real{0.2500}}
  >{\raggedright\arraybackslash}p{(\linewidth - 6\tabcolsep) * \real{0.2500}}@{}}
\toprule\noalign{}
\begin{minipage}[b]{\linewidth}\raggedright
结局
\end{minipage} & \begin{minipage}[b]{\linewidth}\raggedright
全样本，M（SD）
\end{minipage} & \begin{minipage}[b]{\linewidth}\raggedright
开发集，M（SD）
\end{minipage} & \begin{minipage}[b]{\linewidth}\raggedright
留出集，M（SD）
\end{minipage} \\
\midrule\noalign{}
消极情绪 & 0.367 (0.160) & 0.364 (0.159) & 0.388 (0.169) \\
积极情绪 & 0.543 (0.156) & 0.545 (0.154) & 0.524 (0.169) \\
生活满意度 & 0.419 (0.210) & 0.421 (0.211) & 0.401 (0.205) \\
\bottomrule
\end{tabular}
\par\vspace{4pt}注. 完整样本 n = 1,037；开发集 n = 933；留出测试集 n = 104。分数采用所提供 CSV 中的单位，主分析三个结局的范围均为 0--1。每名参与者保留一条最早观察。SD 按样本标准差计算。
\end{minipage}\par\medskip

共享随机森林的开发集宏平均归一化 RMSE 最低（0.9963），独立随机森林紧随其后（0.9965）；训练均值基线为 1.0016。选定的共享森林使用 256 棵树，最小叶节点样本数为 15，每次划分均将全部预测变量作为候选。完整开发集结果和选定设置见补充材料。

\subsection{留出测试集关联与预测误差}\label{ux7559ux51faux6d4bux8bd5ux96c6ux5173ux8054ux4e0eux9884ux6d4bux8befux5dee}

在 104 名留出参与者上，选定共享森林与消极情绪、积极情绪及生活满意度的相关依次为 0.247、0.284 和 0.152（表 2）。经 Holm 校正，检验结果依次为 P = 0.023、P = 0.010 和 P = 0.124。因此，在这一校正下，两个情绪估计可与零区分，而生活满意度的区间包含零。这一模式并不是对不同结局之间差异的检验。图 2 将这些估计与所有已评估的模型类别一并呈现。

\par\medskip\noindent\begin{minipage}{\linewidth}\singlespacing\fontsize{10}{12}\selectfont\setlength{\parindent}{0pt}
{\bfseries 表 2 选定模型与基线的留出测试集表现}\par\vspace{5pt}
\begin{tabular}{@{}
  >{\raggedright\arraybackslash}p{(\linewidth - 8\tabcolsep) * \real{0.1700}}
  >{\raggedright\arraybackslash}p{(\linewidth - 8\tabcolsep) * \real{0.1900}}
  >{\raggedright\arraybackslash}p{(\linewidth - 8\tabcolsep) * \real{0.2500}}
  >{\raggedright\arraybackslash}p{(\linewidth - 8\tabcolsep) * \real{0.1000}}
  >{\raggedright\arraybackslash}p{(\linewidth - 8\tabcolsep) * \real{0.2900}}@{}}
\toprule\noalign{}
\begin{minipage}[b]{\linewidth}\raggedright
模型
\end{minipage} & \begin{minipage}[b]{\linewidth}\raggedright
结局
\end{minipage} & \begin{minipage}[b]{\linewidth}\raggedright
Pearson r［95\% CI］
\end{minipage} & \begin{minipage}[b]{\linewidth}\raggedright
p
\end{minipage} & \begin{minipage}[b]{\linewidth}\raggedright
RMSE［95\% CI］
\end{minipage} \\
\midrule\noalign{}
共享 RF & 消极情绪 & 0.247 {[}0.057, 0.419{]} & 0.012 & 0.166 {[}0.146, 0.186{]} \\
共享 RF & 积极情绪 & 0.284 {[}0.097, 0.452{]} & 0.003 & 0.165 {[}0.143, 0.187{]} \\
共享 RF & 生活满意度 & 0.152 {[}\mbox{-0.042}, 0.335{]} & 0.124 & 0.202 {[}0.181, 0.222{]} \\
训练均值 & 消极情绪 & 无定义 & --- & 0.169 {[}0.150, 0.190{]} \\
训练均值 & 积极情绪 & 无定义 & --- & 0.170 {[}0.146, 0.192{]} \\
训练均值 & 生活满意度 & 无定义 & --- & 0.205 {[}0.184, 0.225{]} \\
独立 RF & 消极情绪 & 0.230 {[}0.039, 0.405{]} & 0.019 & 0.165 {[}0.144, 0.185{]} \\
独立 RF & 积极情绪 & 0.325 {[}0.141, 0.487{]} & \textless{} .001 & 0.163 {[}0.141, 0.186{]} \\
独立 RF & 生活满意度 & 0.184 {[}\mbox{-0.008}, 0.364{]} & 0.061 & 0.201 {[}0.180, 0.222{]} \\
\bottomrule
\end{tabular}
\par\vspace{4pt}注. n = 104。RF 表示随机森林。共享 RF 是开发阶段选定的整体模型，独立 RF 是选定的独立集成模型。Pearson 区间使用 Fisher z；p 值来自未经校正的双侧检验。正文对选定模型的三个相关另作了 Holm 校正。RMSE 区间采用 5,000 次参与者 bootstrap 重抽样；分数保留 CSV 单位。均值预测的 r 无定义。全部八类模型见补充表 S1。
\end{minipage}\par\medskip

\begin{figure}[!htbp]\centering
\includegraphics[width=\linewidth,height=0.54\textheight,keepaspectratio]{figures/figure2_model_performance.pdf}
\par\vspace{5pt}{\raggedright\singlespacing\fontsize{10}{12}\selectfont\noindent 图 2. 八类模型的留出测试集预测表现（n = 104）。A--C：Pearson r 及基于 Fisher-z 的 95\% 区间，采用共同的相关系数坐标范围。D--F：存储分数单位下的 RMSE，以及基于 5,000 次参与者 bootstrap 重抽样的 95\% 区间。菱形和星号标记开发阶段选定的共享 RF，而非统计显著性；空心方块和参考线表示均值基线，其相关系数无定义。模型按开发集表现排序。RF 表示随机森林，ET 表示极端随机树，HGB 表示直方图梯度提升。\par}
\end{figure}

\FloatBarrier

以存储分数单位计，选定模型的 RMSE 依次为 0.166、0.165 和 0.202，决定系数依次为 0.020、0.040 和 0.015。尽管观察值的范围较宽，预测值仍集中在分数中部附近（图 3）。均值基线的 RMSE 依次为 0.169、0.170 和 0.205，表明观察到的数值误差降幅较小。

\begin{figure}[!htbp]\centering
\includegraphics[width=\linewidth,height=0.54\textheight,keepaspectratio]{figures/figure3_observed_predicted.pdf}
\par\vspace{5pt}{\raggedright\singlespacing\fontsize{10}{12}\selectfont\noindent 图 3. 开发阶段选定的共享随机森林所产生的观察分数与预测分数。每个点代表 104 名留出参与者中的一人；两轴采用存储分数单位，并涵盖全部预测值。A：消极情绪；B：积极情绪；C：生活满意度。虚线为表示完全一致的恒等线，而非拟合的回归线。\par}
\end{figure}

\FloatBarrier

\subsection{与基线及独立模型的比较}\label{ux4e0eux57faux7ebfux53caux72ecux7acbux6a21ux578bux7684ux6bd4ux8f83}

消极情绪、积极情绪和生活满意度的``选定模型减去基线''RMSE 差值依次为 \mbox{-0.0034}、\mbox{-0.0048} 和 \mbox{-0.0024}（表 3）。两个情绪差值的逐项 bootstrap 区间均不包含零。然而，没有一项配对交换检验达到六项检验经 Holm 校正后的判定标准（校正后 P = 0.113、P = 0.101 和 P = 1.000）。这些统计量反映不同的不确定性处理方式；多重比较校正并未改变观察到的误差降幅。

``共享随机森林减去独立随机森林''的 RMSE 差值依次为 0.0009、0.0015 和 0.0011。共享程序在每个结局上的误差点估计均略高，但所有差值区间都包含零。这一比较既未确立共享学习的优势，也未证明两类程序等效。模型选择仍然以开发集表现为依据。

\par\medskip\noindent\begin{minipage}{\linewidth}\singlespacing\fontsize{10}{12}\selectfont\setlength{\parindent}{0pt}
{\bfseries 表 3 留出预测误差的配对差异}\par\vspace{5pt}
\begin{tabular}{@{}
  >{\raggedright\arraybackslash}p{(\linewidth - 10\tabcolsep) * \real{0.2200}}
  >{\raggedright\arraybackslash}p{(\linewidth - 10\tabcolsep) * \real{0.1900}}
  >{\raggedright\arraybackslash}p{(\linewidth - 10\tabcolsep) * \real{0.1300}}
  >{\raggedright\arraybackslash}p{(\linewidth - 10\tabcolsep) * \real{0.2200}}
  >{\raggedright\arraybackslash}p{(\linewidth - 10\tabcolsep) * \real{0.1200}}
  >{\raggedright\arraybackslash}p{(\linewidth - 10\tabcolsep) * \real{0.1200}}@{}}
\toprule\noalign{}
\begin{minipage}[b]{\linewidth}\raggedright
比较
\end{minipage} & \begin{minipage}[b]{\linewidth}\raggedright
结局
\end{minipage} & \begin{minipage}[b]{\linewidth}\raggedright
RMSE 差值
\end{minipage} & \begin{minipage}[b]{\linewidth}\raggedright
95\% CI
\end{minipage} & \begin{minipage}[b]{\linewidth}\raggedright
p（交换）
\end{minipage} & \begin{minipage}[b]{\linewidth}\raggedright
p（Holm）
\end{minipage} \\
\midrule\noalign{}
选定模型 vs 均值 & 消极情绪 & \mbox{-0.0034} & {[}\mbox{-0.006}, \mbox{-0.001}{]} & 0.023 & 0.113 \\
选定模型 vs 均值 & 积极情绪 & \mbox{-0.0048} & {[}\mbox{-0.008}, \mbox{-0.001}{]} & 0.017 & 0.101 \\
选定模型 vs 均值 & 生活满意度 & \mbox{-0.0024} & {[}\mbox{-0.007}, 0.002{]} & 0.337 & 1.000 \\
共享 vs 独立 & 消极情绪 & 0.0009 & {[}\mbox{-0.002}, 0.004{]} & 0.533 & 1.000 \\
共享 vs 独立 & 积极情绪 & 0.0015 & {[}\mbox{-0.001}, 0.004{]} & 0.193 & 0.772 \\
共享 vs 独立 & 生活满意度 & 0.0011 & {[}\mbox{-0.002}, 0.004{]} & 0.502 & 1.000 \\
\bottomrule
\end{tabular}
\par\vspace{4pt}注. n = 104。差值为前一个模型减去后一个模型；负值有利于前者。选定／共享模型为共享 RF，独立比较模型为独立 RF。逐项区间采用 5,000 次配对参与者 bootstrap 重抽样。检验通过 10,000 次预测值交换考察均方误差差异，Holm 校正覆盖全部六行。推断以已拟合模型和参与者划分为条件。
\end{minipage}\par\medskip

\subsection{语言类别的贡献}\label{ux8bedux8a00ux7c7bux522bux7684ux8d21ux732e}

消极情绪的平均绝对训练归因值最高的特征为 Sexual（0.0041）、C10（0.0040）和 C2（0.0036）；积极情绪对应为 Sexual（0.0045）、We（0.0039）和 Family（0.0030）；生活满意度对应为 Family（0.0061）、We（0.0056）和 C10（0.0052）（图 4）。这些数值量化了特征对拟合预测的贡献幅度，单位为存储分数单位。C10 和 C2 仍是未映射代码，因而不能依据这些数据赋予其心理学含义。

\begin{figure}[!htbp]\centering
\includegraphics[width=\linewidth,height=0.54\textheight,keepaspectratio]{figures/figure4_feature_attribution.pdf}
\par\vspace{5pt}{\raggedright\singlespacing\fontsize{10}{12}\selectfont\noindent 图 4. 共享 RF 对各结局的十个最大平均绝对训练 TreeSHAP 贡献（n = 933）。A：消极情绪；B：积极情绪；C：生活满意度。数值采用存储结局单位和树路径依赖的背景定义。† 符号与斜线纹理标识尚未核实语义映射的词典代码。绝对值表示贡献大小，不表示方向、统计显著性或因果效应。\par}
\end{figure}

\FloatBarrier

对开发阶段选出的 14 个特征进行留出测试集置换诊断后，误差变化较小且方向不一。一些训练归因较高的特征即使被打乱，也未导致误差恶化。这些结果区分了特征在拟合训练模型中的突出程度与其对留出预测误差的一致影响；完整的逐特征误差变化和洗牌变异见补充材料。

\subsection{预定敏感性分析}\label{ux9884ux5b9aux654fux611fux6027ux5206ux6790}

使用最新记录时，原先选定模型的相关依次为 0.259、0.286 和 0.101。仅保留单次出现 ID 的分析包含 85 名留出参与者，相关依次为 0.303、0.301 和 0.173。排除 ind 和未映射代码后，相关依次为 0.301、0.268 和 0.125，RMSE 为 0.165、0.165 和 0.203。点估计存在变化，但总体上仍呈现关联程度有限、数值预测改善有限的模式。这些敏感性分析针对的是预先指定的数据处理选择；它们既不证明各数据集等效，也不能消除未测量的误差来源。完整区间和开发集重复交叉验证诊断见补充材料。

\subsection{探索性优化}\label{ux63a2ux7d22ux6027ux4f18ux5316}

新增搜索选择了独立森林，其最小叶节点样本量为 60，每次划分考虑一半预测特征。在嵌套开发集验证中，选定流程、固定的原共享森林和均值基线的宏平均归一化 RMSE 分别为 1.0006、1.0001 和 1.0007。选定流程的误差比原森林高 0.044\%，未达到预先规定的探索性判断标准，即降低 1\%。在此前已查看的留出集上，相关分别为 0.245、0.332 和 0.164，RMSE 分别为 0.1661、0.1645 和 0.2023。积极情绪相对均值基线的探索性比较经 Holm 校正后 P = 0.011，而所有结局与原共享森林的比较均未得出明确差异（校正后 P 均为 1.000）。这些重复使用留出集的比较不能确立经独立重复验证的改善。完整结果见补充表 S3--S6；原主分析估计继续保留。

\FloatBarrier

\section{讨论}\label{ux8ba8ux8bba}

\subsection{主要发现}\label{ux4e3bux8981ux53d1ux73b0}

聚合微博语言包含一定程度的参与者积极与消极情绪信息，但估计生活满意度的证据尚不充分。选定模型解释的留出分数变异较少，相对于训练均值的小幅误差降低，在对预定比较进行校正后未得到确立。选定的共享程序也未显示出相对于独立森林的已确立优势。这些发现区分了与自报幸福感之间可检测的关联，以及足够准确的个体评估。由于未直接检验三个结局之间的可预测性差异，结果也不能确立这种差异具有统计显著性。

应结合主观幸福感研究中情感体验与认知评价的区别，解释不同成分的结果 (Diener, 1984)。PANAS 旨在分别测量积极情绪和消极情绪，而 SWLS 要求受访者对整体生活进行评价 (Diener et al., 1985; Watson et al., 1988)。对当前模式的一种可能解释是，相较于生活满意度所依赖的标准与判断，聚合语言类别更容易表征与情绪有关的表达。这是尚未检验的解释，而非相关结果所证明的机制。效标测量、语言覆盖范围或抽样不确定性的差异也可能产生影响。因此，这些结果支持按结局分别开展验证，但尚未说明各结局的数值表现为何不同。

\subsection{与既有研究的关系}\label{ux4e0eux65e2ux6709ux7814ux7a76ux7684ux5173ux7cfb}

这些结果把既有社交媒体语言研究提出的问题，从心理信息是否存在，推进到某种特定表征究竟保留了多少信息。Schwartz 等（2013）发现，开放词汇的 Facebook 特征包含预测信息；Eichstaedt 等（2021）则说明，开放词汇与封闭词汇可以揭示不同类型的模式。这些研究支持在同一样本中比较表征，而不是预先假定准确率一定提高。探索性扩展中的核方法和正则化调整仍使用相同的聚合输入，因此有限的改善并未检验更丰富的语言信息能否提供帮助。Kern 等（2016）还表明，语言数量与聚合方式也会影响分析。由于原始帖子及相应的词级覆盖信息不可得，我们无法区分语言信息不足、效标噪声、抽样波动与模型不匹配的影响。

既有研究也说明，针对单一问卷效标的预测只是评估论证的一部分。Park et al. (2015) 通过聚合与区分关系、知情者报告、外部相关指标和重复语言时间窗评估了基于语言的人格分数。本研究的证据范围更窄：它涉及三个问卷衍生分数的留出估计，尚未确立时间信度或超出这些结局的聚合证据。时间匹配还构成另一项差别。Eichstaedt et al. (2018) 将语言限定为医疗记录首次记载抑郁之前的时期。本数据的最新帖子元数据通常晚于问卷完成时间，且缺少原始帖子，无法设置更早的截止点。因此，当前结果支持回顾性关联，而前瞻性幸福感监测的可能性仍未得到检验。这一比较针对的是验证设计，并不意味着抑郁记录、人格和主观幸福感之间具有等价关系。

其他数字行为研究进一步说明了验证单位与时间安排的重要性。Jaidka 等（2020）发现，地区和社会经济层面的语言差异会干扰基于词典的美国县级幸福感估计。县级估计不能为另一平台上的个体预测提供预期准确率。Stachl 等（2020）使用多类智能手机行为和嵌套验证预测人格，展示了另一种输入与效标组合。De Choudhury 等（2013）明确按照回顾性报告的抑郁发作时间或问卷日期定义 Twitter 文本时间窗；这种设计既不等于前瞻性临床验证，也不能验证连续的幸福感分数。这些比较指出了需要检验的设计选择，而不把不同构念或抽样单位视为可以互换。

共享森林与独立森林的比较，为联合建模幸福感的依据增加了适用条件。多目标树可以通过共同划分降低多个目标的变异，并返回预测向量 (Kocev et al., 2007)。当各目标与预测变量的关系相容时，这种约束可能有用，但仅凭结局间相关并不能确立这种相容性。在本样本中，共享森林在开发集上的微小优势并未转化为已确立的留出优势。与各成分有关的语言关联可能仅部分重叠，也可能是现有数据不足以将小幅收益与抽样变异区分开来。这两种解释均未被直接检验。独立模型为各目标分别选择参数，而共享模型使用一组共同设置，因此这里比较的是选定的程序，而不是对共享结构进行单独操纵。结果不证明等效、负迁移，也不说明多任务学习总体上较差。

\subsection{对幸福感评估的启示}\label{ux5bf9ux5e78ux798fux611fux8bc4ux4f30ux7684ux542fux793a}

这些发现区分了检测语言---幸福感关联与建立评估量尺。模型可能捕捉参与者之间的线性变化，却对准确估计分数贡献有限；本研究中非零的情绪相关、集中于均值附近的预测以及较小的误差降幅，体现了这一点。对预测值进行正比例缩放或平移，可以保持 Pearson 相关不变，却改变 RMSE。这些测量问题属于 Yarkoni and Westfall (2017) 所强调的预测目标与解释目标这一更广泛区分的一部分。可检测的语言---幸福感关联本身并不能识别心理过程、验证所用问卷译本，或确立对个体内变化的敏感性。因此，评估应将拟提出的测量主张与其所需证据联系起来，而不应把显著相关视为一般性的验证结论。

原量表也提出了比单一相关更严格的比较要求。Pavot 和 Diener（1993）区分生活满意度的时间稳定性与对有意义变化的敏感性；对语言推断分数而言，两者都需要重新建立证据。Crawford 和 Henry（2004）在英国成年人样本中评估 PANAS 的心理测量性质，并明确考虑不同群体的因子结构差异。这些证据不能为本研究使用的中文形式或分数变换提供认证。更广泛地说，成熟自评测量的证据包括信度与跨方法聚合关系（Diener et al., 2018）。预测问卷分数，只是为这一验证论证提供一种关系，而不是继承问卷的全部证据基础。经验预测与因果解释之间的区分，也意味着不能把模型准确性解释为已经证实的心理机制（Shmueli, 2010）。

归因结果提供了值得进一步考察的候选特征，但并未确立对幸福感的心理学解释。Family 和 We 等类别与 Sexual 及未映射代码一起，出现在较突出的输入特征中。中文词典验证研究表明，检测到与某个类别相关的词语，并不等同于识别单条帖子的情境性心理含义 (Zhao et al., 2016)。此外，这里报告的绝对 SHAP 幅度，表示特征在所选背景定义下对拟合预测的贡献大小，而不指定真实世界干预的方向或效应 (Janzing et al., 2020; Lundberg et al., 2020)。留出测试集置换引起的误差变化较小且方向不一，也进一步限制了对排序的实质性解释。下一步应核实这些特征的词汇和语境，并检验其贡献能否重复出现，而不是推断谈论某类话题会改善或损害幸福感。

在实践层面，这些模型仍属于研究性估计：它们是否能与直接评估配合使用，需要围绕明确用途进行验证，不能从当前相关结果自动推出。现有误差尚不能换算为原量表上具有实际意义的差异，个体决策效用和总体代表性也都未被评估。任何未来应用还需要保障同意、隐私及人们对敏感推断的控制权。既有社交媒体健康预测研究已明确提出这些问题 (Eichstaedt et al., 2018)，但相关决定必须立足实际数据与预期用途，而不能照搬其他研究的保障性陈述。预测是否可行与是否获准部署，是两个不同问题。

\subsection{局限与未来研究}\label{ux5c40ux9650ux4e0eux672aux6765ux7814ux7a76}

现有数据的若干特征限制了测量方面的结论。实际施测题目措辞、情绪报告时间窗及标签变换均未得到独立核实。原始量表说明能够界定预期构念，却不能确立这些特定存储分数的信度或等价性。评估问卷测量需要题目作答；重建问卷前语言时间窗以及检验奇偶文本一致性，需要带有时间索引的帖子。要判断估计是否能追踪个体内变化，还需要在时间上匹配的重复问卷和相互区分的语言时间窗。本文件中的重复记录并未确立这种匹配；本次参与者间预测也无法回答这一问题。这些缺失输入限定了当前结果的范围，而不是证明相关验证必然失败。

泛化能力也仍有待确定。本分析仅使用一个便利样本和包含 104 名参与者的留出测试集，且没有人口学信息可用于亚组评估。按参与者划分数据并在各折内拟合预处理，可以应对已知的评估泄漏途径 (Kapoor \& Narayanan, 2023)，但并不能确立模型对其他人群、时期或平台的适用性。同样，记录选择和特征排除敏感性分析考察的是特定分析选择，不能替代外部验证。因此，应在独立招募的样本中评估固定的预测器，预先规定可接受误差的标准，并适当覆盖分数拟用于解释的人群。

探索性扩展进一步体现了这种不确定性，而未将其消除。小样本可能产生不稳定的性能估计，交叉验证训练集之间的重叠也会使简单假设折间独立的计算产生误导（Varoquaux, 2018）。因此，我们保留逐折结果，不把折间标准差当作置信区间，也不把重复使用的留出集上有利的比较解释为独立确认。新增搜索未呈现稳定的开发集优势。这削弱了继续调整同一特征矩阵模型的理由，但不能据此认定社交媒体幸福感预测普遍存在某个准确率上限。

在明确分数定义、词典和时间匹配后，两项针对性比较尤其有价值。首先，在相同参与者划分下比较固定类别与语境特征或开放词汇特征，可以检验表征是否限制了各个幸福感成分的估计。其次，以足够精度和针对各结局的诊断重复共享程序与独立程序的比较，可以将可重复的共享学习收益与微小、不稳定的差异区分开来。两项比较均不应仅以相关是否跨越显著性阈值来判断。真正相关的进展，是证明所得分数能针对明确用途改善误差、稳定性或解释，同时继续以参与者报告的体验为依据。

\section{结论}\label{ux7ed3ux8bba}

在留出参与者中，微博心理语言特征与情绪分数呈现一定程度的线性关联，相较均值基线的观察误差降幅较小。生活满意度的估计尚不充分，且未直接检验结局之间的差异。指定的共享输出程序未显示出优于独立森林的优势。这些发现为现有表征和学习程序提供了评估基准；要建立幸福感评估工具，还需要关于数值准确性、信度、时间匹配和泛化能力的进一步证据。

\section{声明}\label{ux58f0ux660e}

\subsection{数据与代码可用性}\label{ux6570ux636eux4e0eux4ee3ux7801ux53efux7528ux6027}

数据集与分析产物保存在本地研究包中，尚未公开存储，也未建立持久标识符。数据保管方、共享权限、访问条件及数据集引用均待作者补充。稳定的参与者 ID 不应被视为已经实现不可逆匿名化。

\subsection{伦理与作者声明}\label{ux4f26ux7406ux4e0eux4f5cux8005ux58f0ux660e}

伦理批准和知情同意信息据原始采集说明报告，并需由作者确认其对此次二次分析的适用性。作者身份、贡献、经费、利益冲突和致谢均待作者在独立标题页中补充。

\subsection{人工智能的使用}\label{ux4ebaux5de5ux667aux80fdux7684ux4f7fux7528}

OpenAI Codex 协助完成了分析代码开发与执行、文献检索和证据核查、图表制作及写作。修订还使用了结构化科学写作指导和基于问题的一致性审查。自动检查及独立代理审查不构成负责的人类作者的认可；对数据来源、方法、解释、参考文献和声明的审核仍待完成。工具与技能来源记录在修订档案中。

\setstretch{1.08}\fontsize{10.5}{12.8}\selectfont

\FloatBarrier

\section{参考文献}\label{ux53c2ux8003ux6587ux732e}

\APAref Breiman, L. (2001). Random forests. \emph{Machine Learning}, \emph{45}(1), 5--32. \url{https://doi.org/10.1023/A:1010933404324}

\APAref Caruana, R. (1997). Multitask learning. \emph{Machine Learning}, \emph{28}(1), 41--75. \url{https://doi.org/10.1023/A:1007379606734}

\APAref Cawley, G. C., \& Talbot, N. L. C. (2010). On over-fitting in model selection and subsequent selection bias in performance evaluation. \emph{Journal of Machine Learning Research}, \emph{11}(70), 2079--2107. \url{https://jmlr.org/papers/v11/cawley10a.html}

\APAref Crawford, J. R., \& Henry, J. D. (2004). The Positive and Negative Affect Schedule (PANAS): Construct validity, measurement properties and normative data in a large non-clinical sample. \emph{British Journal of Clinical Psychology}, \emph{43}(3), 245--265. \url{https://doi.org/10.1348/0144665031752934}

\APAref De Choudhury, M., Gamon, M., Counts, S., \& Horvitz, E. (2013). Predicting depression via social media. \emph{Proceedings of the International AAAI Conference on Web and Social Media}, \emph{7}(1), 128--137. \url{https://doi.org/10.1609/icwsm.v7i1.14432}

\APAref Diener, E. (1984). Subjective well-being. \emph{Psychological Bulletin}, \emph{95}(3), 542--575. \url{https://doi.org/10.1037/0033-2909.95.3.542}

\APAref Diener, E., Emmons, R. A., Larsen, R. J., \& Griffin, S. (1985). The satisfaction with life scale. \emph{Journal of Personality Assessment}, \emph{49}(1), 71--75. \url{https://doi.org/10.1207/s15327752jpa4901_13}

\APAref Diener, E., Suh, E. M., Lucas, R. E., \& Smith, H. L. (1999). Subjective well-being: Three decades of progress. \emph{Psychological Bulletin}, \emph{125}(2), 276--302. \url{https://doi.org/10.1037/0033-2909.125.2.276}

\APAref Diener, E., Lucas, R. E., \& Oishi, S. (2018). Advances and open questions in the science of subjective well-being. \emph{Collabra: Psychology}, \emph{4}(1), 15. \url{https://doi.org/10.1525/collabra.115}

\APAref Eichstaedt, J. C., Smith, R. J., Merchant, R. M., Ungar, L. H., Crutchley, P., Preoţiuc-Pietro, D., Asch, D. A., \& Schwartz, H. A. (2018). Facebook language predicts depression in medical records. \emph{Proceedings of the National Academy of Sciences}, \emph{115}(44), 11203--11208. \url{https://doi.org/10.1073/pnas.1802331115}

\APAref Eichstaedt, J. C., Kern, M. L., Yaden, D. B., Schwartz, H. A., Giorgi, S., Park, G., Hagan, C. A., Tobolsky, V. A., Smith, L. K., Buffone, A., Iwry, J., Seligman, M. E. P., \& Ungar, L. H. (2021). Closed- and open-vocabulary approaches to text analysis: A review, quantitative comparison, and recommendations. \emph{Psychological Methods}, \emph{26}(4), 398--427. \url{https://doi.org/10.1037/met0000349}

\APAref Geurts, P., Ernst, D., \& Wehenkel, L. (2006). Extremely randomized trees. \emph{Machine Learning}, \emph{63}(1), 3--42. \url{https://doi.org/10.1007/s10994-006-6226-1}

\APAref Jaidka, K., Giorgi, S., Schwartz, H. A., Kern, M. L., Ungar, L. H., \& Eichstaedt, J. C. (2020). Estimating geographic subjective well-being from Twitter: A comparison of dictionary and data-driven language methods. \emph{Proceedings of the National Academy of Sciences}, \emph{117}(19), 10165--10171. \url{https://doi.org/10.1073/pnas.1906364117}

\APAref Janzing, D., Minorics, L., \& Blöbaum, P. (2020). Feature relevance quantification in explainable AI: A causal problem. In S. Chiappa \& R. Calandra (Eds.), \emph{Proceedings of the Twenty Third International Conference on Artificial Intelligence and Statistics} (Vol. 108, pp. 2907--2916). PMLR. \url{https://proceedings.mlr.press/v108/janzing20a.html}

\APAref Kapoor, S., \& Narayanan, A. (2023). Leakage and the reproducibility crisis in machine-learning-based science. \emph{Patterns}, \emph{4}(9), 100804. \url{https://doi.org/10.1016/j.patter.2023.100804}

\APAref Kern, M. L., Park, G., Eichstaedt, J. C., Schwartz, H. A., Sap, M., Smith, L. K., \& Ungar, L. H. (2016). Gaining insights from social media language: Methodologies and challenges. \emph{Psychological Methods}, \emph{21}(4), 507--525. \url{https://doi.org/10.1037/met0000091}

\APAref Kocev, D., Vens, C., Struyf, J., \& Džeroski, S. (2007). Ensembles of multi-objective decision trees. In J. N. Kok, J. Koronacki, R. López de Mántaras, S. Matwin, D. Mladenič, \& A. Skowron (Eds.), \emph{Machine Learning: ECML 2007} (Vol. 4701, pp. 624--631). Springer Berlin Heidelberg. \url{https://doi.org/10.1007/978-3-540-74958-5_61}

\APAref Lundberg, S. M., Erion, G., Chen, H., DeGrave, A., Prutkin, J. M., Nair, B., Katz, R., Himmelfarb, J., Bansal, N., \& Lee, S.-I. (2020). From local explanations to global understanding with explainable AI for trees. \emph{Nature Machine Intelligence}, \emph{2}(1), 56--67. \url{https://doi.org/10.1038/s42256-019-0138-9}

\APAref Park, G., Schwartz, H. A., Eichstaedt, J. C., Kern, M. L., Kosinski, M., Stillwell, D. J., Ungar, L. H., \& Seligman, M. E. P. (2015). Automatic personality assessment through social media language. \emph{Journal of Personality and Social Psychology}, \emph{108}(6), 934--952. \url{https://doi.org/10.1037/pspp0000020}

\APAref Pavot, W., \& Diener, E. (1993). Review of the Satisfaction With Life Scale. \emph{Psychological Assessment}, \emph{5}(2), 164--172. \url{https://doi.org/10.1037/1040-3590.5.2.164}

\APAref Ryan, R. M., \& Deci, E. L. (2001). On happiness and human potentials: A review of research on hedonic and eudaimonic well-being. \emph{Annual Review of Psychology}, \emph{52}(1), 141--166. \url{https://doi.org/10.1146/annurev.psych.52.1.141}

\APAref Schwartz, H. A., Eichstaedt, J. C., Kern, M. L., Dziurzynski, L., Ramones, S. M., Agrawal, M., Shah, A., Kosinski, M., Stillwell, D., Seligman, M. E. P., \& Ungar, L. H. (2013). Personality, gender, and age in the language of social media: The open-vocabulary approach. \emph{PLoS ONE}, \emph{8}(9), e73791. \url{https://doi.org/10.1371/journal.pone.0073791}

\APAref Shmueli, G. (2010). To explain or to predict? \emph{Statistical Science}, \emph{25}(3), 289--310. \url{https://doi.org/10.1214/10-STS330}

\APAref Stachl, C., Au, Q., Schoedel, R., Gosling, S. D., Harari, G. M., Buschek, D., Völkel, S. T., Schuwerk, T., Oldemeier, M., Ullmann, T., Hussmann, H., Bischl, B., \& Bühner, M. (2020). Predicting personality from patterns of behavior collected with smartphones. \emph{Proceedings of the National Academy of Sciences}, \emph{117}(30), 17680--17687. \url{https://doi.org/10.1073/pnas.1920484117}

\APAref Tausczik, Y. R., \& Pennebaker, J. W. (2010). The psychological meaning of words: LIWC and computerized text analysis methods. \emph{Journal of Language and Social Psychology}, \emph{29}(1), 24--54. \url{https://doi.org/10.1177/0261927X09351676}

\APAref Varoquaux, G. (2018). Cross-validation failure: Small sample sizes lead to large error bars. \emph{NeuroImage}, \emph{180}, 68--77. \url{https://doi.org/10.1016/j.neuroimage.2017.06.061}

\APAref Watson, D., Clark, L. A., \& Tellegen, A. (1988). Development and validation of brief measures of positive and negative affect: The PANAS scales. \emph{Journal of Personality and Social Psychology}, \emph{54}(6), 1063--1070. \url{https://doi.org/10.1037/0022-3514.54.6.1063}

\APAref Yarkoni, T., \& Westfall, J. (2017). Choosing prediction over explanation in psychology: Lessons from machine learning. \emph{Perspectives on Psychological Science}, \emph{12}(6), 1100--1122. \url{https://doi.org/10.1177/1745691617693393}

\APAref Zhang, Y., \& Yang, Q. (2022). A survey on multi-task learning. \emph{IEEE Transactions on Knowledge and Data Engineering}, \emph{34}(12), 5586--5609. \url{https://doi.org/10.1109/TKDE.2021.3070203}

\APAref Zhao, N., Jiao, D., Bai, S., \& Zhu, T. (2016). Evaluating the validity of simplified Chinese version of LIWC in detecting psychological expressions in short texts on social network services. \emph{PLOS ONE}, \emph{11}(6), e0157947. \url{https://doi.org/10.1371/journal.pone.0157947}

\end{document}
